\section{Síntesis y ampliación}

Los aprendizajes clave de esta sesión:

\begin{enumerate}
  \item El \textbf{ancho equivalente} $k \times m = 6144$ es la clave para entender el costo computacional del MoE: equivale a un FFN Dense de ancho $3 \times d_{\text{model}}$
  \item La suma ponderada y la concatenación son equivalentes en FLOPs y en el rank de la matriz; el MoE aporta diversidad mediante varios Expert «delgados»
  \item La atención es la interacción causal entre tokens, mientras que el FFN/MoE es un mapeo independiente dentro de cada token: los papeles de ambos son, en esencia, distintos
  \item Ya sea Dense o MoE, la parte del FFN sigue siendo, con mucho, la que concentra la mayor cantidad de parámetros
  \item RMSNorm aporta 2 por capa más 1 al final; el número de parámetros es pequeño, pero conceptualmente no se puede pasar por alto
\end{enumerate}

\subsection{Lecturas adicionales}

\begin{itemize}
  \item Código fuente de HuggingFace Transformers: \texttt{models/qwen3\_moe/}
  \item Switch Transformer (Fedus et al., 2021): una aplicación clásica del MoE en el Transformer
  \item Informe técnico de DeepSeek MoE: un diseño de Expert de grano más fino
\end{itemize}

\end{document}
